\begin{theorem}[The common core on a trace]
  \label{thm:aba-gather-commoncore}
  \lean{PLTS.ABA.Gather.CoreTrace}\lean{PLTS.ABA.Gather.specInst_core}\lean{PLTS.ABA.Gather.instanceOverBroadcastSpecification_core}\lean{PLTS.ABA.Gather.instanceOverBracha_core}\leanok
  \uses{def:aba-gather-spec, def:aba-gather-over-broadcast-specification, def:aba-gather-over-bracha, thm:gather-core, thm:gather-broadcast-substitution, thm:trace-support}
  $\mathrm{CoreTrace}$ is the instance-level reading of the core, quantified over the labels of a run rather than over the states of one: every announcement $\mathrm{core}(C)$ of the trace names a set $C$ of at least $n - f$ entries, any two announcements of the trace name the same set, and a return $\mathrm{ret}(id, g)$ and an announcement $\mathrm{core}(C)$ of the trace satisfy $C \sqsubseteq g$.
  Every trace in the support of every achievable trace distribution of the gather specification instance satisfies it.\lean{PLTS.ABA.Gather.specInst_core}\leanok{} So does every such trace of the implementation over broadcast specifications and of the implementation over Bracha instances.\lean{PLTS.ABA.Gather.instanceOverBroadcastSpecification_core}\lean{PLTS.ABA.Gather.instanceOverBracha_core}\leanok{} This is Transition System~4's binding clause read on a trace.
\end{theorem}

\begin{proof}\leanok
  \uses{def:aba-gather-spec, thm:gather-core, thm:gather-broadcast-substitution, thm:trace-support, thm:forward-sound}
  At the specification the three clauses come from the transitions alone.
  \leandeclnamed{bindCore}{PLTS.ABA.Gather.Step.bindCore} is the only writer of $core$ and fires only from $core = \bot$, so the field keeps its value once written and every announcement of an execution names the same set; the size bound is that transition's own guard, carried forward as a state invariant by induction along the execution; and a return requires $C \sqsubseteq g$ of the core $C$ standing at its state, which is the set every announcement of the execution names.
  Theorem~\ref{thm:trace-support} turns a positive-probability trace into a genuine execution whose return labels and announcements are steps of those transitions, which is the trace form.
  The two implementations inherit the predicate along the soundness of their refinements: the implementation over broadcast specifications along Theorem~\ref{thm:gather-core}, and the implementation over Bracha instances along Theorem~\ref{thm:gather-broadcast-substitution} composed with it.
  What makes that transfer say anything is that the core is announced on a label of its own.
  The implementations hold it in a field no process reads, the refinements match the announcements, and a trace-support property transfers along an inclusion of achievable trace distributions with no hypothesis beyond positive probability.
\end{proof}
